\section{Structural Core and Closure Principles}
\label{sec:structural-core}
%% ============================================================

Before compressing the preceding results into a single ToE master equation, we
make explicit the structural kernel that has been implicit throughout the
construction.

\begin{definition}[Structural core]
\label{def:structural-core}
The structural core of the theory consists of the following coupled
requirements:
\begin{enumerate}
\item a common underlying field object $\Psi$;
\item a triolic decomposition
\[
  \Psi = \Psi_1\oplus\Psi_2\oplus\Psi_3;
\]
\item an orthogonal closure/light sector $\Psi_3$ relative to the first two
sectors;
\item a dual unitary symmetry exchanging $\Psi_1$ and $\Psi_2$;
\item a Heisenberg compensator $C$ emerging from the ground condition and
selecting the collapse/light output;
\item a projected closure condition determining the physically admissible
sector;
\item a phase-shifted observer frame in which effective quantities are read
out.
\end{enumerate}
\end{definition}

\begin{theorem}[Existence of a nontrivial closure-stable triolic system]
\label{thm:structural-existence}
Assume the structural core of Definition~\ref{def:structural-core}. Then there
exists a nontrivial class of configurations in which the triolic decomposition,
compensator consistency, and projected closure condition are simultaneously
satisfied.
\end{theorem}

\begin{proof}
The preceding sections already exhibit the required ingredients separately.
The central theorem provides a common ground object together with its triolic
unfolding and Lorentz-compatible strata. The unitary symmetry $J$ links the
first two sectors, while the Heisenberg compensator produces the
collapse/light sector as a structurally distinguished output. The subsequent
probe--closure analysis introduces the projected closure residue and identifies
an admissible closed sector by the condition that leakage into the complement
is suppressed. Since these ingredients are realized within one and the same
construction, the intersection of the triolic, compensator-consistent, and
closure-stable constraints is nonempty. Hence a nontrivial closure-stable
triolic system exists.
\end{proof}

\begin{proposition}[Centrality of the Heisenberg compensator]
\label{prop:hc-centrality}
Within the structural core, the Heisenberg compensator is not an auxiliary
correction term but the central structural hinge linking duality, collapse,
closure stability, and the phase-shifted observer description.
\end{proposition}

\begin{proof}
Without the duality symmetry, the first two sectors would remain merely
parallel descriptions. Without the compensator, their common collapse/light
output would not be selected. Without this selected output, the projected
closure condition would not isolate a physically admissible sector. And
without that closure-stable sector, the observer-frame description would lack a
well-defined quantity to read out after phase shift. Thus the compensator is
precisely the operator through which duality becomes collapse, collapse becomes
closure, and closure becomes observable effective physics.
\end{proof}


\subsection*{Preliminary Semantic and Tensorial Kernel Interface}
\addcontentsline{toc}{subsection}{Preliminary Semantic and Tensorial Kernel Interface}

This subsection states the semantic/tensor interface in anticipatory form. Its formal unfolding as a lemma--theorem chain is given later in Phase II-E.

The structural kernel can be stated on two coordinated levels. The first is a
semantic layer of sets, states, maps, and admissibility predicates. The second
is a tensorial realization layer in which the same constraints are written in a
notation compatible with the Riemann/Sphaera core. The tensor layer does not
replace the semantic one; it compresses it into a transportable operator form.

\begin{definition}[Semantic structural kernel]
\label{def:semantic-structural-kernel}
The semantic structural kernel consists of the data
\[
(\mathbb S,\{S_n\}_{n\ge 0},S_0,\hat K,J,C,\mathcal M)
\]
together with the admissibility predicates
\[
\begin{aligned}
&\mathrm{Triolic}(x),\ \mathrm{Anchored}(x),\ \mathrm{FrameShifted}(x),\ \mathrm{GroundLinked}(x),\\
&\mathrm{HCForced}(x),\ \mathrm{ClosureStable}(x),\ \mathrm{HigherConfined}(x),\ \mathrm{SeptinionReadable}(x).
\end{aligned}
\]
These objects satisfy the following semantic conditions:
\begin{enumerate}[leftmargin=2em,label=(S\arabic*)]
\item \textbf{Cascade hierarchy:}
\[
\mathbb S=S_N\supseteq S_{N-1}\supseteq \cdots \supseteq S_1\supseteq S_0,
\qquad
\hat K=\Pi_0\circ\Pi_1\circ\cdots\circ\Pi_{N-1}.
\]
\item \textbf{Ground-state terminality:}
\[
x\in S_0 \Longleftrightarrow \hat K(x)=x.
\]
\item \textbf{Maximal-totality closure rule:} if $M^*\cap \mathbb S=\emptyset$, then
\[
\hat K(M^*)\subseteq S_0.
\]
\item \textbf{$J$-involution:}
\[
J^2=\mathrm{id}.
\]
\item \textbf{Compensator map:}
\[
C(x):=\tfrac12\bigl(x+J(x)\bigr).
\]
\item \textbf{Triolic minimality:} if $\mathrm{Triolic}(x)$, then $x$ admits a three-sector
organization $(K_1,K_2,K_3)$.
\item \textbf{Orthogonal anchor:} if $\mathrm{Anchored}(x)$, then the third sector is
orthogonal to the first two,
\[
K_3\perp K_1,
\qquad
K_3\perp K_2.
\]
\item \textbf{Frame-shift principle:} $\mathrm{FrameShifted}(x)$ means that $x$ is a shifted
or regime-dependent readout of a common structural body rather than an
ontologically independent object.
\end{enumerate}
\end{definition}

\begin{definition}[Tensor realization of the structural kernel]
\label{def:tensor-structural-kernel}
A tensor realization of Definition~\ref{def:semantic-structural-kernel}
is given by the operator family
\[
\mathcal J^\mu{}_{\nu},\qquad
\mathcal C^\mu{}_{\nu},\qquad
\Pi_n{}^\mu{}_{\nu},\qquad
\hat K^\mu{}_{\nu},\qquad
\mathcal M^\mu{}_{\nu},
\]
subject to the following identities:
\begin{enumerate}[label=(T\arabic*)]
  \item \textbf{Involution:}
    $\mathcal J^\mu{}_{\alpha}\mathcal J^\alpha{}_{\nu}=\delta^\mu{}_{\nu}$.
  \item \textbf{Compensator idempotence:}
    $\mathcal C^\mu{}_{\alpha}\mathcal C^\alpha{}_{\nu}=\mathcal C^\mu{}_{\nu}$.
  \item \textbf{Cascade composition:}
    $\hat K^\mu{}_{\nu}=(\Pi_0\Pi_1\cdots\Pi_{N-1})^\mu{}_{\nu}$.
  \item \textbf{$\mathcal M$-admissibility (intertwining):}
    The Millennium matrix $\mathcal M^\mu{}_{\nu}$ is required to be
    \emph{admissible} in the sense that it preserves the algebraic
    structure of the kernel:
    \[
      \mathcal M^\mu{}_{\alpha}\mathcal J^\alpha{}_{\nu}
      =\mathcal J^\mu{}_{\alpha}\mathcal M^\alpha{}_{\nu},
      \qquad
      \mathcal M^\mu{}_{\alpha}\mathcal C^\alpha{}_{\nu}
      =\mathcal C^\mu{}_{\alpha}\mathcal M^\alpha{}_{\nu}.
    \]
    Whenever layer transport is required to preserve the cascade reading,
    additionally
    \[
      \mathcal M^\mu{}_{\alpha}\hat K^\alpha{}_{\nu}
      =\hat K^\mu{}_{\alpha}\mathcal M^\alpha{}_{\nu}.
    \]
\end{enumerate}
The tensor $\mathcal J$ represents the involutive dual/frame symmetry,
$\mathcal C$ represents the compensator projection, the $\Pi_n$ represent
cascade projections between adjacent strata, and $\mathcal M$ is the
Millennium matrix as an admissible structure-preserving converter between
admissible layers. Condition~(T4) is the defining admissibility condition
on $\mathcal M$; it is not derivable from (T1)--(T3) alone, but is imposed
as the minimal requirement for $\mathcal M$ to qualify as a
structure-preserving converter in the sense of
Definition~\ref{def:semantic-structural-kernel}.
\end{definition}

\begin{proposition}[Semantic forcing chain]
\label{prop:semantic-forcing-chain}
Within the semantic structural kernel, the following implication chain holds:
\[
\mathrm{GroundLinked}(x)\wedge \mathrm{FrameShifted}(x)
\Longrightarrow \mathrm{HCForced}(x),
\]
\[
\mathrm{Triolic}(x)\wedge \mathrm{Anchored}(x)\wedge \mathrm{HCForced}(x)
\Longrightarrow \mathrm{ClosureStable}(x),
\]
\[
\mathrm{ClosureStable}(x)\wedge \mathrm{HigherConfined}(x)
\Longrightarrow \mathrm{SeptinionReadable}(x).
\]
\end{proposition}

\begin{proof}
\textbf{First implication:
$\mathrm{GroundLinked}(x)\wedge\mathrm{FrameShifted}(x)
\Rightarrow\mathrm{HCForced}(x)$.}

$\mathrm{GroundLinked}(x)$ means $x\in S_n$ for some stratum $n$; it is
reachable by the cascade $\hat{K}$ from Definition~\ref{def:semantic-structural-kernel}~(S1).
$\mathrm{FrameShifted}(x)$ means by~(S8) that $x$ is a shifted or
regime-dependent readout of a common structural body.
By Lemma~\ref{lem:one-body-lemma}, the shifted presentations all refer
to one and the same carrier. By Lemma~\ref{lem:re-identification-lemma},
admissibility then requires a re-identification mechanism returning
competing presentations to a common output channel. The only operator
in the semantic kernel with this role is the compensator $C(x)=\tfrac{1}{2}(x+J(x))$
from~(S5), which is the unique $J$-symmetric projection.
Hence $\mathrm{HCForced}(x)$: the compensator is structurally necessary.

\smallskip
\textbf{Second implication:
$\mathrm{Triolic}(x)\wedge\mathrm{Anchored}(x)\wedge\mathrm{HCForced}(x)
\Rightarrow\mathrm{ClosureStable}(x)$.}

This is Lemma~\ref{lem:closure-forcing-lemma}, now with an explicit
four-step proof citing
Lemmas~\ref{lem:triolic-minimality-lemma},~\ref{lem:orthogonal-anchor-lemma},
\ref{lem:hc-admissibility-lemma},~\ref{lem:hc-idempotence-lemma},
and Theorem~\ref{thm:fractal}.

\smallskip
\textbf{Third implication:
$\mathrm{ClosureStable}(x)\wedge\mathrm{HigherConfined}(x)
\Rightarrow\mathrm{SeptinionReadable}(x)$.}

$\mathrm{ClosureStable}(x)$ means $\hat{K}(\mathcal{C}x)=\mathcal{C}x$
(Lemma~\ref{lem:closure-stable-iff-im-C}):
$x\in\mathrm{im}(C)$ is already in the ground sector.
$\mathrm{HigherConfined}(x)$ means $x$ is contained in a higher coupled
algebraic container — specifically the septinion/seption embedding
(Lemma~\ref{lem:higher-confinement-lemma}).
Within the higher container, the closure-stable state is readable as a
seption element precisely because its compensator-stabilized image
$\mathcal{C}x$ already lies in the admissible ground sector $S_0$, and
the higher container adds no new ontology
(Lemma~\ref{lem:non-escape-lemma}: no admissible transition exports $x$
into an independent external sector).
Hence the higher-confinement reading is a restatement of the already
established closure at higher algebraic level:
$\mathrm{SeptinionReadable}(x)$.
\end{proof}

\begin{proposition}[Tensorial intertwining conditions for the Millennium matrix]
\label{prop:tensorial-intertwining-m}
Every admissible tensor realization in the sense of
Definition~\ref{def:tensor-structural-kernel} satisfies:
\[
\mathcal M\mathcal J=\mathcal J\mathcal M,
\qquad
\mathcal M\mathcal C=\mathcal C\mathcal M,
\]
and, whenever layer transport is required to preserve the cascade reading,
\[
\mathcal M\hat K=\hat K\mathcal M.
\]
\end{proposition}

\begin{proof}
These are precisely condition~(T4) of
Definition~\ref{def:tensor-structural-kernel}, written in operator
notation. A tensor realization is admissible by definition only if
$\mathcal M$ satisfies the intertwining identities~(T4); hence the
relations hold for every admissible realization. No additional argument
is required beyond the definition.

The structural motivation for imposing~(T4) is as follows. The first
relation $\mathcal M\mathcal J=\mathcal J\mathcal M$ ensures that
transport by $\mathcal M$ does not destroy the involutive dual/frame
symmetry encoded by $\mathcal J$~(T1). The second relation
$\mathcal M\mathcal C=\mathcal C\mathcal M$ ensures that the
compensator-selected admissible sector $\mathrm{im}(\mathcal C)$ is
$\mathcal M$-invariant, consistent with $\mathcal C$ being idempotent
(T2). The third relation, when imposed, says that layer conversion
commutes with the cascade readout~(T3). These three motivations
explain why~(T4) is the correct admissibility condition; they do not
constitute an independent proof of~(T4) from (T1)--(T3).
\end{proof}

\begin{remark}[Resolution of Roadmap item A2]
\label{rem:a2-resolved}
Roadmap item~A2 identified that the intertwining relation
$\mathcal M\mathcal C=\mathcal C\mathcal M$ was supported only by a
necessity argument rather than a derivation. The resolution adopted
in v2.27.6 is to promote~(T4) to an explicit definitional condition
on admissible tensor realizations
(Definition~\ref{def:tensor-structural-kernel}, condition~(T4)).
Proposition~\ref{prop:tensorial-intertwining-m} then follows by
definition. The conditional marker on
Theorem~\ref{thm:hc-spectral-projector}
(Remark~\ref{rem:hc-spectral-projector-conditional}) is therefore
resolved: the theorem is no longer conditional on an unproven
hypothesis, but on an explicit definitional requirement.
\end{remark}

\begin{remark}[Semantic--tensor correspondence]
\label{rem:semantic-tensor-correspondence}
The semantic predicates and the tensor notation are intended to be read in
parallel. In particular,
\[
\mathrm{HCForced}(x)
\quad\leftrightarrow\quad
x\mapsto \mathcal Cx,
\]
\[
\mathrm{ClosureStable}(x)
\quad\leftrightarrow\quad
\hat K(\mathcal Cx)=\mathcal Cx,
\]
and
\[
\mathrm{Admissible}(x)
\quad\Rightarrow\quad
\mathrm{Admissible}(\mathcal Mx)
\]
whenever $\mathcal M$ satisfies the intertwining conditions of
Proposition~\ref{prop:tensorial-intertwining-m}. This makes it possible to
write later structure conditions in the same tensorial style as the
Riemann/Sphaera core while keeping their semantic meaning explicit.
\end{remark}


\subsection*{Phase I: First Unfolding Lemmas for the Generative Layer}
\addcontentsline{toc}{subsection}{Phase I: First Unfolding Lemmas for the Generative Layer}

We now begin the actual mathematical unfolding of the normalized structural
core by isolating the first chain of generative lemmas. The aim is not yet to
expand the full field/probe realization, but to make explicit the initial
forcing route from a common structural body to triolic organization,
orthogonal anchoring, frame-shifted readout, re-identification, and finally the
necessity of the Heisenberg compensator.

\begin{lemma}[One-body lemma]
\label{lem:one-body-lemma}
The semantic and tensorial kernels are to be read as describing different
admissible presentations of one and the same structural body, not independent
ontological objects. In particular, the cascade hierarchy, dual/frame symmetry,
and compensator projection all act on a common carrier rather than on unrelated
state spaces.
\end{lemma}

\begin{proof}
In Definition~\ref{def:semantic-structural-kernel} the kernel data are given as
one coordinated family
\[
(\mathbb S,\{S_n\}_{n\ge 0},S_0,\hat K,J,C,\mathcal M),
\]
not as several unrelated structures. The grouped architecture of
Definition~\ref{def:grouped-structural-architecture} likewise states that the
generative, projection, and field/probe layers are different organizational
regimes of the same underlying body. The tensor realization of
Definition~\ref{def:tensor-structural-kernel} then presents the same body in
operator form. Hence the structural reading is one-body throughout.
\end{proof}

\begin{lemma}[Triolic minimality lemma]
\label{lem:triolic-minimality-lemma}
A merely binary decomposition is structurally insufficient for the normalized
core. The minimal nondegenerate organization compatible with admissible
selection, compensator action, and closure readout is triolic.
\end{lemma}

\begin{proof}
By Definition~\ref{def:semantic-structural-kernel}, a triolic state carries a
three-sector organization $(K_1,K_2,K_3)$, while the orthogonal anchor singles
out the third sector as stabilizing direction relative to the first two. If one
retained only a binary split, then duality between the first two sectors could
still be stated, but there would be no distinguished third channel in which the
compensator-selected output and closure/light readout could be structurally
stabilized. The normalized core therefore requires three sectors as the minimal
nondegenerate organization.
\end{proof}

\begin{lemma}[Orthogonal anchor lemma]
\label{lem:orthogonal-anchor-lemma}
Within a triolic organization, the third sector is not an optional decoration but
the structural anchor that prevents the first two sectors from remaining a
merely parallel or relabelable pair. In particular,
\[
K_3\perp K_1,
\qquad
K_3\perp K_2.
\]
\end{lemma}

\begin{proof}
The orthogonal anchor condition is built into
Definition~\ref{def:semantic-structural-kernel}. Its role is clarified by the
forcing chain of Proposition~\ref{prop:semantic-forcing-chain}: closure
stability does not follow from duality alone, but from triolic structure plus a
protected anchor together with HC forcing. The third sector therefore fixes the
otherwise ambiguous two-sector pairing and provides the admissible direction in
which collapse/light selection can be stabilized.
\end{proof}

\begin{lemma}[Frame-shift lemma]
\label{lem:frame-shift-lemma}
Admissible effective states are read out in a phase- or frame-shifted
presentation. Accordingly, an effective state is not introduced as an
independent object but as a shifted view of the common structural body.
\end{lemma}

\begin{proof}
This is exactly the content of the frame-shift principle in
Definition~\ref{def:semantic-structural-kernel}, where
$\mathrm{FrameShifted}(x)$ means that $x$ is a shifted or regime-dependent
readout of a common body. The grouped architecture makes the same point at the
layer level: projection/translation and field/probe realizations are
organizational regimes of one underlying structure and therefore inherit rather
than replace the generative body.
\end{proof}

\begin{lemma}[Re-identification lemma]
\label{lem:re-identification-lemma}
If a state is simultaneously one-body and frame-shifted, then admissibility
requires a re-identification mechanism that returns competing presentations to a
common structural channel.
\end{lemma}

\begin{proof}
By Lemma~\ref{lem:one-body-lemma}, the relevant descriptions concern one and the
same body. By Lemma~\ref{lem:frame-shift-lemma}, these descriptions appear in
shifted or regime-dependent presentations. Without re-identification, the
shifted presentations would remain only formally related and no unique
admissible output sector could be selected. But the semantic kernel is built to
support admissible cascade readout and closure selection; hence a mechanism that
returns the shifted presentations to a common channel is structurally required.
\end{proof}

\begin{theorem}[HC-forcing theorem]
\label{thm:hc-forcing-theorem}
Within the normalized structural core, the required re-identification mechanism
is the Heisenberg compensator. Equivalently,
\[
\mathrm{GroundLinked}(x)\wedge \mathrm{FrameShifted}(x)
\Longrightarrow \mathrm{HCForced}(x).
\]
\end{theorem}

\begin{proof}
Proposition~\ref{prop:semantic-forcing-chain} already states the forcing
implication. The present point is to identify why the implication is necessary.
\[
\begin{aligned}
&\text{ground condition}
\Longrightarrow
\text{one triolic projection}
\Longrightarrow
\text{triolic minimality}
\Longrightarrow
\text{orthogonal anchor}
\\[0.2em]
&\Longrightarrow
\text{frame-shifted readout}
\Longrightarrow
\text{re-identification}
\Longrightarrow
\mathrm{HC}.
\end{aligned}
\]
\end{proof}

\begin{corollary}[First generative forcing chain]
\label{cor:first-generative-forcing-chain}
The first unfolded forcing route of the normalized core is
\[
\begin{aligned}
\text{one body}
&\Longrightarrow \text{triolic minimality}\\
&\Longrightarrow \text{orthogonal anchor}\\
&\Longrightarrow \text{frame-shifted readout}\\
&\Longrightarrow \text{re-identification}\\
&\Longrightarrow \mathrm{HC}.
\end{aligned}
\]
This chain is the entry point for the later closure and higher-confinement
unfolding.
\end{corollary}

\begin{proof}
Lemmas~\ref{lem:one-body-lemma}--\ref{lem:re-identification-lemma} establish the
successive structural requirements, and
Theorem~\ref{thm:hc-forcing-theorem} identifies their first nontrivial forced
operator consequence. The later closure and higher-confinement chain therefore
starts from this generative route.
\end{proof}


\subsection*{Phase I-B: HC and Closure Unfolding}
\addcontentsline{toc}{subsection}{Phase I-B: HC and Closure Unfolding}

We next unfold the first closure layer that begins once HC has been forced.
The aim of this block is to make precise that HC is not merely present, but
acts as an admissibility-preserving, anchor-compatible, closure-selecting
mechanism, and that exact closure can be distinguished from residual or only
partial closure.

\begin{lemma}[HC idempotence lemma]
\label{lem:hc-idempotence-lemma}
Once a state has been compensator-normalized, repeated compensator action does
not move it to a different admissible sector. Equivalently,
\[
C(C(x)) = C(x).
\]
\end{lemma}

\begin{proof}
By Definition~\ref{def:tensor-structural-kernel}, the tensorial compensator
satisfies $\mathcal C^2=\mathcal C$, and by
Definition~\ref{def:semantic-structural-kernel} the semantic compensator is the
corresponding map
\[
C(x)=\tfrac12\bigl(x+J(x)\bigr).
\]
Thus compensator action is projective rather than drifting: once a state lies
in the compensator-selected channel, further application leaves it there.
\end{proof}

\begin{lemma}[HC admissibility lemma]
\label{lem:hc-admissibility-lemma}
If $x$ is structurally admissible, then its compensator-normalized image $C(x)$
remains admissible.
\end{lemma}

\begin{proof}
The semantic/tensor kernel was introduced precisely so that the compensator is
not an external deformation but an internal normalization operator. Since
Theorem~\ref{thm:hc-forcing-theorem} identifies HC as the required
re-identification mechanism, using $C$ restores competing shifted presentations
to a common admissible channel instead of producing a foreign state. In the
tensor realization this is the admissible output selected by $\mathcal C$; on
the semantic side the same role is expressed by the preservation of the common
structural body under compensator projection. Hence admissibility is preserved.
\end{proof}

% ------------------------------------------------------------------
% Spectral-projector block (v2.27.4 integration)
% ------------------------------------------------------------------

\begin{definition}[Spectral decomposition of the state space]
\label{def:spectral-decomposition-M}
Let $\mathcal{M}$ denote the Millennium matrix acting on the admissible state
space $\mathcal{S}$. A \emph{spectral decomposition of $\mathcal{S}$ under
$\mathcal{M}$} is an orthogonal direct sum
\[
  \mathcal{S} \;=\; \bigoplus_{\lambda} \mathcal{S}_\lambda,
\]
where each $\mathcal{S}_\lambda \subset \mathcal{S}$ is a closed
$\mathcal{M}$-invariant subspace satisfying
$\mathcal{M}|_{\mathcal{S}_\lambda} = \lambda\,\mathrm{id}$.
A state $s \in \mathcal{S}$ is called \emph{spectrally stable} if
$s \in \mathcal{S}_\lambda$ for some $\lambda$, i.e.\ $\mathcal{M}(s) =
\lambda\,s$.
\end{definition}

\begin{remark}
Definition~\ref{def:spectral-decomposition-M} extends the existing
decomposition $\Hplus = \Hplusplus \oplus \Hplusminus$
(cf.\ Theorem~\ref{thm:unity}) to the full admissible readout regime.
The two primary spectral subspaces under any $\mathcal{M}$ satisfying
Proposition~\ref{prop:tensorial-intertwining-m} are precisely $\Hplusplus$
and $\Hplusminus$, since both are $\mathcal{M}$-invariant by the intertwining
condition $\mathcal{M}\mathcal{J} = \mathcal{J}\mathcal{M}$.
\end{remark}

\begin{lemma}[Diagonal self-traversal forces scalar action on the compensated sector]
\label{lem:diagonal-traversal-scalar-action}
Assume the Sphaera self-traversal condition of Theorem~\ref{thm:fractal}:
the cascade operator $\hat{K}$ transports the null axis $L_n$ of each
stratum $S_n$ to the null axis $L_{n-1}$ of the preceding stratum,
carrying the same structural content without generating new independent
directions. Then the Millennium matrix $\mathcal{M}$, when restricted to
the image of the compensator $\mathrm{im}(\mathcal{C})$, acts as a scalar
multiple of the identity:
\[
  \mathcal{M}\big|_{\mathrm{im}(\mathcal{C})} \;=\; \lambda_0\,\mathrm{id},
\]
where $\lambda_0$ is the unique admissible eigenvalue of $\mathcal{M}$ on
the closure-stable sector.
\end{lemma}

\begin{proof}
By Theorem~\ref{thm:fractal}, the self-traversal of the Sphaera is
\emph{diagonal}: the cascade does not generate a new structural object at
each stratum, but transports one and the same null-axis content through
admissible realization levels. In operator language this means that the
cascade direction at each stratum is not a Jordan-chain descendant of the
previous one---no new vector is generated from an existing vector by a
nilpotent shift.

By the intertwining condition $\mathcal{M}\mathcal{C} = \mathcal{C}\mathcal{M}$
(Proposition~\ref{prop:tensorial-intertwining-m}), $\mathrm{im}(\mathcal{C})$
is $\mathcal{M}$-invariant. Suppose for contradiction that
$\mathcal{M}|_{\mathrm{im}(\mathcal{C})}$ had a non-trivial Jordan block:
then there would exist $s,w \in \mathrm{im}(\mathcal{C})$ with $w\neq 0$ and
\[
  \mathcal{M}(s) = \lambda_0\,s + w.
\]
This would mean $\mathcal{M}$ generates $w$ from $s$---a new independent
direction within the compensated sector---contradicting the diagonal
self-traversal property. Hence $\mathcal{M}|_{\mathrm{im}(\mathcal{C})}$
is diagonalizable. Since $\mathrm{im}(\mathcal{C})$ is the closure-stable
sector (Lemma~\ref{lem:closure-forcing-lemma}), which is irreducible under
$\mathcal{M}$, a diagonalizable action on an irreducible invariant subspace
is scalar. Denoting this scalar $\lambda_0$ completes the proof.
\end{proof}

\begin{theorem}[HC as spectral projector of the Millennium matrix]
\label{thm:hc-spectral-projector}
Assume the admissible state space $\mathcal{S}$ carries a spectral
decomposition under $\mathcal{M}$ as in
Definition~\ref{def:spectral-decomposition-M}, that $\mathcal{M}$ satisfies
the intertwining conditions of Proposition~\ref{prop:tensorial-intertwining-m},
and that the Sphaera self-traversal is diagonal in the sense of
Lemma~\ref{lem:diagonal-traversal-scalar-action}. Then the Heisenberg
compensator $C$ coincides with the spectral projector of $\mathcal{M}$
onto the closure-stable eigenspace $\mathcal{S}_0 :=
\ker(\mathcal{M} - \lambda_0\,\mathrm{id})$:
\[
  C \;=\; \Pi^{(\mathcal{M})}_{\mathrm{spec}}.
\]
\end{theorem}

\begin{remark}[Status of Theorem~\ref{thm:hc-spectral-projector}]
\label{rem:hc-spectral-projector-conditional}
As of v2.27.6, both open premises of this theorem are resolved:
the bijection $\mathrm{ClosureStable}(x)\Leftrightarrow x\in\mathrm{im}(C)$
is established by Lemma~\ref{lem:closure-stable-iff-im-C} (Roadmap A3),
and the intertwining relation $\mathcal{M}\mathcal{C}=\mathcal{C}\mathcal{M}$
is now a definitional condition (T4) of
Definition~\ref{def:tensor-structural-kernel} (Roadmap A2,
Remark~\ref{rem:a2-resolved}).
The theorem therefore holds unconditionally within every admissible
tensor realization.
\end{remark}

\begin{proof}
We verify the three conditions characterizing a spectral projector.

\smallskip
\noindent\textbf{Step 1: Idempotence.}
By Lemma~\ref{lem:hc-idempotence-lemma}, $C^2=C$. By
Definition~\ref{def:tensor-structural-kernel},
$\mathcal{C}^\mu{}_\alpha\mathcal{C}^\alpha{}_\nu = \mathcal{C}^\mu{}_\nu$
in the tensor realization. Hence $C$ is a projector.

\smallskip
\noindent\textbf{Step 2: $\mathcal{M}$-compatibility.}
By Proposition~\ref{prop:tensorial-intertwining-m},
$\mathcal{M}\mathcal{C} = \mathcal{C}\mathcal{M}$. Hence $\mathrm{im}(C)$
is $\mathcal{M}$-invariant and $C$ projects onto a $\mathcal{M}$-stable
subspace.

\smallskip
\noindent\textbf{Step 3: Scalar eigenvalue on the image.}
By Lemma~\ref{lem:diagonal-traversal-scalar-action}, the diagonal
self-traversal forces $\mathcal{M}|_{\mathrm{im}(C)} = \lambda_0\,\mathrm{id}$.
Hence $\mathrm{im}(C) \subseteq \mathcal{S}_0$.
By Lemma~\ref{lem:closure-forcing-lemma}, the closure-stable states are
exactly those in $\mathrm{im}(C)$, so $\mathrm{im}(C) = \mathcal{S}_0$.

\smallskip
\noindent\textbf{Conclusion.}
A projector with $\mathcal{M}$-invariant image equal to a single eigenspace
$\mathcal{S}_0$ is the spectral projector $\Pi^{(\mathcal{M})}_{\mathrm{spec}}$
onto that eigenspace.
\end{proof}

\begin{corollary}[Structural fixpoint: $\mathcal{M}$ and $C$ mutually determine each other]
\label{cor:M-HC-fixpoint}
Under the hypotheses of Theorem~\ref{thm:hc-spectral-projector},
\[
  \mathcal{M}\,C \;=\; C\,\mathcal{M} \;=\; \lambda_0\,C.
\]
\end{corollary}

\begin{proof}
For any $s\in\mathcal{S}$: since $C(s)\in\mathcal{S}_0$, we have
$\mathcal{M}(C(s))=\lambda_0\,C(s)$, giving $\mathcal{M}\,C=\lambda_0\,C$.
For $s\in\mathcal{S}_0$: $C\,\mathcal{M}(s)=C(\lambda_0 s)=\lambda_0 C(s)$.
For $s\in\ker(C)$: $C\,\mathcal{M}(s)=0$ (since $\ker(C)$ is
$\mathcal{M}$-invariant and maps to a sector annihilated by $C$).
Hence $C\,\mathcal{M}=\lambda_0\,C$.
\end{proof}

\begin{remark}[Phase-shifted frame and spectral measurement]
\label{rem:frame-shift-spectral}
The identification $C = \Pi^{(\mathcal{M})}_{\mathrm{spec}}$ has a direct
consequence for the matter frame. Since matter observers reside at
$\zminus = +\bsym \approx 0.5493$ rather than the photon ground shell
$\zminus = 0$, every spectral readout of $\mathcal{M}$ is performed in the
shifted frame. The measured eigenvalue is therefore not $\lambda_0$ itself
but the frame-corrected value
\[
  \lambda_0^{(\mathrm{matter})}
  \;=\; \gL^{-k(\lambda_0)}\cdot\lambda_0^{(\mathrm{photon})},
\]
where $k(\lambda_0)$ is the frame-insertion count of the observable and
$\gL \approx 1.864$ is the matter--photon Lorentz factor
(Theorem~\ref{thm:frame-offset}). All physical constants derived from
spectral eigenvalues of $\mathcal{M}$ therefore carry systematic
$\gL^k$-corrections relative to their photon-frame values.
\end{remark}

\begin{corollary}[Standard Model realization of the spectral decomposition]
\label{cor:SM-spectral-realization}
The octonionic decomposition $\OO = \mathbb{C}\oplus\mathbb{C}^3$
(Theorem~\ref{thm:octonion-SU3}) is a concrete realization of
Definition~\ref{def:spectral-decomposition-M}: the singlet $\mathbb{C}$ is
the closure-stable eigenspace $\mathcal{S}_0 = \mathrm{im}(C)$, and the
triplet $\mathbb{C}^3$ carries the three color eigenspaces
$\mathcal{S}_{K_1},\mathcal{S}_{K_2},\mathcal{S}_{K_3}$.
Under this identification, Theorem~\ref{thm:hc-spectral-projector}
specializes as follows:
\begin{enumerate}[label=(\roman*)]
  \item The compensator $C$ projects onto the singlet $\mathbb{C}\subset\OO$,
  which is the photon kernel $K_3$ sector ($\zminus=0$).
  \item The triplet $\mathbb{C}^3$ is $\mathcal{M}$-invariant and carries
  the SU(3) color action; matter ($K_1$) and tachyon ($K_2$) sectors are
  spectral subspaces at non-zero frame offset.
  \item The matter-frame correction reads
  \[
    \lambda_0^{(\mathrm{matter})}
    = \gL^{-k(\lambda_0)}\cdot\lambda_0^{(\mathrm{photon})},
  \]
  consistent with Theorem~\ref{thm:SM}: $K_3$ is frame-invariant
  ($k=0$), while the triplet components carry $k\ge 1$ corrections
  and are therefore observed with $\gL^k$-suppression in the matter frame.
\end{enumerate}
\end{corollary}

\begin{proof}
By Theorem~\ref{thm:octonion-SU3}, the octonion automorphism group $G_2$
restricts to $\mathrm{SU}(3)\subset G_2$, yielding the direct sum
$\OO = \mathbb{C}\oplus\mathbb{C}^3$. The singlet $\mathbb{C}$ is the
unique $\mathrm{SU}(3)$-invariant subspace, hence $\mathcal{M}$-stable at
a scalar eigenvalue $\lambda_0$. By the remark in
Theorem~\ref{thm:SM}, $K_3 = e_1 e_2$ is simultaneously the photon kernel
and the projection of the Heisenberg compensator; hence
$\mathrm{im}(C) = \mathbb{C} = \mathcal{S}_0$. The three kernels
$(K_1,K_2,K_3)$ realize the three colors of $\mathbb{C}^3$, confirming
$\mathcal{S}_{K_i}$ as spectral subspaces. Statement~(iii) then follows
directly from Remark~\ref{rem:frame-shift-spectral} applied to each
triplet component.
\end{proof}

% ------------------------------------------------------------------
% Readout-Relative Lemmas and Structural Lifts (v2.27.4 integration)
% ------------------------------------------------------------------

